\section{CONCLUSÃO}\label{sec:conclusao}

A principal conclusão deste artigo é que a dependência de transferências, como a contabilidade a registra, não acompanha a autonomia fiscal. O indicador ordena corretamente os municípios, mas erra o nível de quem tem a receita vinculada à própria economia, e pode cair enquanto a autonomia cai. O que importa é quem controla a receita e de onde ela vem, não o rótulo que ela recebe no orçamento, e é por isso que índices de autonomia mais completos medem essas dimensões em separado \citep{Ladner2016}.

Em Camaçari, boa parte da ``dependência'' é a própria economia do município voltando ao caixa da prefeitura pela cota do ICMS. No papel, mais da metade de sua receita vem de transferências; sem a cota, menos de um terço. O município não controla essa base, mas também não depende de decisões alheias para recebê-la: está vinculado à própria economia, para o bem e para o mal, como mostrou a saída da Ford. Barreiras, ao contrário, depende de fato de decisões tomadas fora, porque o FPM segue a arrecadação federal. Essa vinculação tem um custo distributivo: ao devolver o ICMS a quem produz mais, a regra concentra recursos onde já há riqueza \citep{daCosta2024}. Vinculação à base e redistribuição puxam em sentidos opostos, e a reforma tributária, ao repartir o futuro IBS sobretudo pela população \citep{EC132}, escolhe a redistribuição: tende a desligar a receita dos municípios industriais de sua própria economia e a dar mais recursos aos demais.

A Hungria mostra o caminho inverso, com menos dependência e menos autonomia: pelos instrumentos descritos na Seção~\ref{sec:hungria_lit}, o governo central reduziu o controle local sobre funções e receitas sem que isso aparecesse como aumento do peso das transferências. Não há paradoxo nisso, porque funções e recursos saíram juntos; há apenas um indicador que não foi feito para captar controle.

Quanto ao efeito \textit{flypaper}, o gasto dos dois municípios acompanha as transferências muito mais do que o PIB, mas acompanha na mesma medida o FUNDEB, e com dois municípios não é possível separar o efeito da regra de gastar o que se arrecada. Assim, H1 fica sem confirmação; H2 vale para a dependência total, mas não para as transferências incondicionais; H3 é rejeitada; e H4 recebe apoio descritivo nos dois países.

O estudo contribui com uma base de dados aberta e corrigida, que revelou, entre outros pontos, resultados orçamentários negativos em Barreiras antes registrados como superávits, e com um alerta para quem avalia finanças municipais: indicadores de autonomia precisam separar a receita que nasce da economia local da que vem de fora. Com dois municípios, as conclusões não se generalizam nem identificam causas. Um painel com os municípios do interior na faixa máxima do FPM \citep{Bernardelli2026, Castro2021221} e dados municipais sobre as transferências discricionárias na Hungria são os próximos passos naturais.
